\documentclass[11pt]{article}
\usepackage[a4paper,margin=24mm]{geometry}
\usepackage{amsmath,amssymb,booktabs,longtable,hyperref}
\hypersetup{hidelinks}
\setlength{\emergencystretch}{2em}
\newcommand{\Q}{\mathbb Q}
\newcommand{\val}{\operatorname{ord}}
\newcommand{\den}{\operatorname{den}}
\newcommand{\lcm}{\operatorname{lcm}}
\title{Arithmetic and combinatorics of the pure-$X$ coefficients\\An exact scan through genus ten}
\author{}\date{16 October 2026}
\begin{document}\maketitle
\noindent
The scan covers 144 pure-$X$ coefficients and 16,384 nonzero monomial
coefficients of the physical quintic polynomials $P_2,\ldots,P_{10}$.
The count 16,384 is simply the observed total; no significance is assigned
to its being a power of two. Genera nine and ten were computed after the
first patterns had been identified in genera two through eight. All arithmetic
and comparisons below are exact. The leading-edge formulas are derived
consequences of the imposed BCOV and conifold equations. The broader denominator
bound remains a conjectural interpretation of the finite data.

\section{What is being scanned}
Retain the original YY conventions of the implementation:
\[
 P_g=C^{g-1}F_g=P_g^{\rm an}+a_g(X),\qquad
 a_g(X)=\sum_{n=0}^{d_g}a_{g,n}X^n,\qquad d_g=3g-3,
\]
where $P_g^{\rm an}(v_1=v_2=v_3=0,X)=0$.
The coefficient of $X^n$ in the full polynomial is generally a polynomial
in $v_1,v_2,v_3$. The phrase ``pure-$X$ coefficient'' here means specifically
$[u^0v_1^0v_2^0v_3^0X^n]P_g$.
For denominator tests, the scan also examines all the mixed monomials.

Set $x=z/\mu=\psi^{-5}$, $\delta=1-x$, and $X=1-\delta^{-1}$.
The rational coefficients in this normalized ring do not depend on $\mu$.
A coordinate conversion back to $z$ introduces known external powers of
$\mu$. Consequently keeping $\mu$ free removes coordinate-scale clutter, but
does not itself create or alter the arithmetic patterns below.
All statements are specific to the stated generator and amplitude normalization.

\section{The highest pure-$X$ coefficient: an exact formula}
Define $b_g=|B_{2g}|$. The result is
\begin{equation}\boxed{
 A_g:=a_{g,3g-3}=5^{3g-3}\frac{b_g}{2g(2g-2)}.}
 \label{eq:leading}
\end{equation}
The normalized values are in Table~\ref{tab:edge}. The apparently exceptional
numerators $691$ at genus six and $3617$ at genus eight are precisely the
Bernoulli numerators, not new unexplained integers introduced by the solver.

\begin{table}[ht]\centering
\begingroup\renewcommand{\arraystretch}{1.35}
\begin{tabular}{rrl}
\toprule $g$ & $3g-3$ & $a_{g,3g-3}/5^{3g-3}$ \\
\midrule
2 & 3 & $\frac{1}{240}$ \\
3 & 6 & $\frac{1}{1008}$ \\
4 & 9 & $\frac{1}{1440}$ \\
5 & 12 & $\frac{1}{1056}$ \\
6 & 15 & $\frac{691}{327600}$ \\
7 & 18 & $\frac{1}{144}$ \\
8 & 21 & $\frac{3617}{114240}$ \\
9 & 24 & $\frac{43867}{229824}$ \\
10 & 27 & $\frac{174611}{118800}$ \\
\bottomrule
\end{tabular}
\endgroup

\caption{Highest pure-$X$ coefficient with the known power of five removed.
Every row agrees exactly with the independently computed physical polynomial.}
\label{tab:edge}
\end{table}

Here is the derivation, including the signs. In the rational conifold frame,
\[
 f_c=1+O(\delta^3),\quad
 \tau=\nu/f_c=\delta+\tfrac7{10}\delta^2+O(\delta^3),\quad
 K_g=\left[-\frac{\delta f_c^2}{625}\right]^{g-1}.
\]
The physical amplitude is $\mathcal F_g^c=K_gP_g^{\rm hol}$, and the gap is
\[
 \mathcal F_g^c=\kappa_g\tau^{2-2g}+O(1),\qquad
 \kappa_g=\frac{5^{1-g}b_g}{2g(2g-2)}.
\]
The generators $v_i$ are regular in this frame. Weighted degree implies that
only $A_gX^{d_g}$ contributes to the most singular pole.
Since $X^{d_g}=(-1)^{d_g}\delta^{-d_g}(1-\delta)^{d_g}$, its sign after
multiplication by $K_g$ is $(-1)^{d_g+g-1}=(-1)^{4g-4}=+1$.
Thus $A_g/625^{g-1}=\kappa_g$, which gives \eqref{eq:leading}.
This is a consequence of the gap normalization already imposed, rather than
an independent new boundary condition.

\subsection{Determination from lower-order data}
Even the Bernoulli input has an exact convolution in lower indices:
\begin{equation}\boxed{
 b_1=\frac16,\qquad
 b_g=\frac1{2g+1}\sum_{h=1}^{g-1}\binom{2g}{2h}b_hb_{g-h},\quad g\ge2.}
 \label{eq:bernoulli_rec}
\end{equation}
This is the classical Bernoulli convolution with its alternating signs removed.
One way to derive it is to use
\[
 H(t)=\frac t2\cot\frac t2
     =1-\sum_{g\ge1}\frac{b_g}{(2g)!}t^{2g},\qquad
 tH'=H-H^2-\frac{t^2}{4}.
\]
At degree $2g\ge4$, equating coefficients gives
$(2g+1)b_g/(2g)!=\sum_{h=1}^{g-1}b_hb_{g-h}/((2h)!(2g-2h)!)$.
The scan checks this convolution against a separate linear Bernoulli recurrence
through $g=53$. Formula \eqref{eq:leading} therefore predicts this edge of the
polynomial through genus 53 without computing full genus-53 amplitudes.

\section{A second exact law: the next pure-$X$ coefficient}
Put
\[
 k_h=5^{3h-2}\frac{b_h}{2h},\quad h\ge1;
 \qquad k_1=\frac5{12}.
\]
Then
\begin{equation}\boxed{\begin{aligned}
 a_{g,3g-4}={}&-\frac85(g-1)A_g
       +5(2g-3)k_{g-1}
       +\sum_{h=1}^{g-1}k_hk_{g-h}.
\end{aligned}}\label{eq:next}\end{equation}
This is an actual genus-splitting convolution. It matches every genus
$2,\ldots,10$, including the fresh genus-nine and genus-ten computations.
For example it gives $a_{2,2}=205/144$, $a_{3,5}=33125/1008$,
and $a_{4,8}=18453125/12096$.
It also predicts the sign change in $a_{g,3g-4}$ between genera four and five.

To derive it, let $P_h^{(1)}$ be the normalized one-point vertex and define
\[
 M_g=[v_1X^{3g-4}]P_g.
\]
At $u=v_1=v_2=v_3=0$, the highest pure-$X$ term of $P_h^{(1)}$ is
$k_hX^{3h-2}$. For $h\ge2$, differentiating $A_hX^{3h-3}$ gives
$5(3h-3)A_hX^{3h-2}$; the connection contributes
$-5(h-1)A_hX^{3h-2}$. Their sum is $10(h-1)A_h=k_h$.
The same formula includes the separately known $h=1$ seed $5X/12$.
A second covariant derivative has leading pure-$X$ coefficient
$5(2h-1)k_h$. Therefore the $u^0$ anomaly equation gives
\begin{equation}
 M_g=-\frac12\left[5(2g-3)k_{g-1}
                       +\sum_{h=1}^{g-1}k_hk_{g-h}\right].
\end{equation}
This mixed coefficient is also checked directly against every full polynomial.

At the conifold $v_1(0)=2$. Weighted degree shows that the next pole of $P_g$
can receive contributions only from $A_gX^{d_g}$,
$a_{g,d_g-1}X^{d_g-1}$, and $M_gv_1X^{d_g-1}$.
The next term in $K_gP_g$ is consequently
\[
 625^{1-g}\bigl[-d_gA_g-a_{g,d_g-1}-2M_g\bigr]\delta^{3-2g}.
\]
Meanwhile $\tau^{2-2g}=\delta^{2-2g}
[1-\frac75(g-1)\delta+O(\delta^2)]$.
Equating these coefficients gives
$a_{g,d_g-1}=-\frac85(g-1)A_g-2M_g$, hence \eqref{eq:next}.
This derivation uses only the leading coefficients of earlier genera and the
specified local frame, without fitting a sequence to the data.

\section{A triangular binomial rule determines a larger block}
There is a useful extension beyond these two explicit formulas. Given the
known anomaly body at genus $g$, form
\[
 R_g(\delta)=\frac{\kappa_g\tau^{2-2g}}{K_g}
             -P_g^{\rm an}(v_1^c,v_2^c,v_3^c,X_c),\qquad
 r_{g,j}=[\delta^{-j}]R_g(\delta).
\]
The regular $O(1)$ part of the physical amplitude, divided by $K_g$, is
$O(\delta^{1-g})$. Thus $r_{g,j}$ is prescribed for every $j\ge g$.
But $X=1-\delta^{-1}$ implies
\[
 [\delta^{-j}]a_g(1-\delta^{-1})
 =(-1)^j\sum_{n=j}^{3g-3}\binom nj a_{g,n}.
\]
Consequently there is an exact downward recursion
\begin{equation}\boxed{
 a_{g,j}=(-1)^j r_{g,j}
       -\sum_{n=j+1}^{3g-3}\binom nj a_{g,n},
 \quad j=3g-3,\ldots,g.}
\end{equation}
It determines the highest $2g-2$ pure-$X$ coefficients. Unlike the special
formula for the next coefficient, the general rule uses the full lower-genus
amplitudes through the anomaly body, not just their pure-$X$ subsequences.
An independent implementation of this triangular rule predicts and checks
90 coefficients across genera two through ten, without using those unknown
coefficients in constructing the predictions.
The remaining low powers still require the other boundary data.

\section{The constant term is also controlled}
The MUM constant-map condition fixes
\begin{equation}\boxed{
 a_{g,0}=100\,625^{g-1}
 \frac{|B_{2g}B_{2g-2}|}{2g(2g-2)(2g-2)!}.}
\end{equation}
It is positive in this normalization. This supplies another exactly
predictable coefficient, with a product of Bernoulli numbers and a factorial.
It is useful to keep this known contribution visible when searching for
arithmetic novelty in denominator tables.

\section{A proved sieve and an empirical one}
\subsection{Bernoulli denominators: an established theorem}
The von Staudt--Clausen theorem gives
\[
 \den(B_{2g})=\prod_{\substack{p\text{ prime}\\p-1\mid2g}}p.
\]
Together with \eqref{eq:leading}, this completely explains the denominator
of the leading coefficient. If $B_{2g}=N_g/D_g$ is reduced, then
\[
 \den(A_g)=\frac{2g(2g-2)D_g}
 {\gcd\bigl(5^{3g-3}|N_g|,\,2g(2g-2)D_g\bigr)}.
\]
The exact numerator still contains Bernoulli arithmetic, but it is generated
by \eqref{eq:bernoulli_rec}. This is a precise sense in which both
combinatorics and a prime-divisibility sieve control the same coefficients.
When studying denominators, dividing out $5^{3g-3}$ is transparent.
Dividing by the entire leading coefficient would insert its Bernoulli numerator
into other denominators and can create misleading appearances of primes such
as $691$.

\subsection{Full coefficient arrays: observed Todd-type bound}
Let $L_g$ be the least common multiple of all reduced coefficient denominators
of $P_g$. Let $L_g^{\ge7}$ retain only its prime factors $p\ge7$.
For every computed genus $2\le g\le10$, we find the exact equality
\begin{equation}
 L_g^{\ge7}=\prod_{\substack{p\ge7\\p\text{ prime}}}
                    p^{\lfloor2g/(p-1)\rfloor}.
 \label{eq:empirical}
\end{equation}
The pure-$X$ coefficient LCM has the same $p\ge7$ part.
Every coefficient is also 5-integral: its denominator is coprime to 5.
No denominator prime larger than $2g+1$ occurs in the data.

The product in \eqref{eq:empirical} is the $p\ge7$ part of the
Todd/Hirzebruch denominator of degree $2g$.
The Hirzebruch denominator formula is a theorem for Todd polynomials
\cite{BV}; its applicability as a bound to all our BCOV coefficients
has not been proved here. The natural conjecture suggested by the data is
\[
 \val_p(L_g)\le\left\lfloor\frac{2g}{p-1}\right\rfloor,
 \qquad p\ge7,
\]
with equality throughout the observed range. We do not claim eventual equality
for every genus, nor a proof of this upper bound.

\begin{table}[ht]\centering
\begin{tabular}{rrrrl}
\toprule $g$ & Terms & $\val_2(L_g)$ & $\val_3(L_g)$ & $L_g^{\ge7}$ \\
\midrule
2 & 14 & 5 & 2 & $1$ \\
3 & 62 & 7 & 4 & $7$ \\
4 & 185 & 11 & 5 & $7$ \\
5 & 435 & 13 & 6 & $7\cdot11$ \\
6 & 877 & 16 & 8 & $7^{2}\cdot11\cdot13$ \\
7 & 1590 & 18 & 9 & $7^{2}\cdot11\cdot13$ \\
8 & 2666 & 23 & 10 & $7^{2}\cdot11\cdot13\cdot17$ \\
9 & 4211 & 25 & 13 & $7^{3}\cdot11\cdot13\cdot17\cdot19$ \\
10 & 6344 & 28 & 14 & $7^{3}\cdot11^{2}\cdot13\cdot17\cdot19$ \\
\bottomrule
\end{tabular}

\caption{Exact denominators in the full YY polynomials. The factors at 2 and 3
are displayed separately; the last column is the observed $p\ge7$ part.}
\end{table}

There are two qualifications to the apparent strength of this observation.
First, primes 2 and 3 have larger exponents than the naive Todd degree-$2g$
formula in some genera, while powers of 5 are absent. The generator
normalization and genus-one seed matter.
Second, throughout $2\le g\le10$, the single known constant-map coefficient
$a_{g,0}$ already attains all the maximal exponents at primes $p\ge7$.
Thus the new empirical information in the rest of the array is the upper bound:
none of the other 16,000-plus coefficients introduces a larger exponent.
The observed LCM equality is not an independent determination of their numerators.

The bound would still be useful. Polynomial primitives can divide by exponents
as large as $3g-3$; the absence of primes between $2g+1$ and $3g-3$ therefore
requires cancellations. For example, at genus ten the degree bound permits
27, but no denominator factor 23 appears.
A proof would have to explain those cancellations in the normalized recursion,
perhaps through a stable-graph or characteristic-class organization. This is a
research direction suggested by the data, not an established identification.

A weaker additional observation is
\[
 \min_{\text{nonzero coefficients }c\text{ of }P_g}\val_5(c)=
 \begin{cases}1,&2\le g\le6,\\2,&7\le g\le10.\end{cases}
\]
The staircase $\lceil(g-1)/5\rceil$ fits these values, but only one jump has
been seen. It is recorded as a low-confidence extrapolation, not promoted
to a law. It predicts the next increase at genus twelve.

\section{Searches that did not give a predictive law}
The scan also tested first-order rational-ratio recurrences
\[
 a_{j+1}/a_j=P(j)/Q(j),\qquad\deg P,\deg Q\le2.
\]
They were tested both across genus along the top four diagonals (after
removing $5^{3g-3}$) and across $X$ powers within each available row.
Fits used the first five consecutive ratios. Every fitted model failed at
held-out coefficients. In particular, a degree-$(2,2)$ rational function can
interpolate five ratios, so fitting those five data points alone is not evidence.
This is only a bounded search; it does not exclude other recurrence structures.

We have not found a compact new rule fixing all interior numerators from
nearby pure-$X$ coefficients. They remain controlled by the full anomaly and
boundary system. The evidence supports separating the known Bernoulli and
binomial structures from an integer numerator array before attempting further
recurrence or congruence searches.

\section{Reproduction and precise scope}
Run from the extracted directory:
\begin{verbatim}
python scan_coefficients.py physical/amplitudes.json --output results
python verify_triangular.py
\end{verbatim}
The first script uses only the Python standard library. The second uses the
bundled snapshot of the exact algebra and period modules, also without external
runtime dependencies. Input data, per-genus physical diagnostics for the newly
computed genera, and the complete exact results are included.
The file \texttt{exact\_edge\_predictions\_through\_53.json} contains
formulas \eqref{eq:leading} and \eqref{eq:next}, the constant term, and the
mixed coefficient $M_g$ through genus 53. It does not contain full amplitudes
above genus ten. The scan checks 45 edge/arithmetic identities against the
nine physical polynomials, the Bernoulli convolution through index $g=53$,
and 90 independent triangular coefficient predictions.

\begin{thebibliography}{9}
\bibitem{YY} S. Yamaguchi and S.-T. Yau,
\emph{Topological String Partition Functions as Polynomials},
\href{https://arxiv.org/abs/hep-th/0406078}{hep-th/0406078},
especially the polynomial anomaly equations and the stable-graph discussion.
\bibitem{BV} V. M. Buchstaber and A. P. Veselov,
\emph{Todd polynomials and Hirzebruch numbers},
\href{https://arxiv.org/abs/2310.07383}{arXiv:2310.07383v4}.
This paper discusses both the Hirzebruch denominator formula and the
von Staudt--Clausen formula, together with combinatorial representations.
\end{thebibliography}
\end{document}
